
% ===== Slide 47 =====
\section{Solution Sets}
\begin{frame}{Solution Sets of Linear Systems}
  \centering
  \Large 1.5 Solution Sets of Linear Systems
\end{frame}

% ===== Slide 48 =====
\begin{frame}{Topics}
  \begin{itemize}
    \item Parametric Vector Form
    \item Solution Sets of Homogeneous Linear Systems
    \item Solution Sets of Nonhomogeneous Systems
  \end{itemize}
\end{frame}

% ===== Slide 49 =====
\begin{frame}{Parametric Vector Form}
  A parametric description
  \[
    \begin{cases}
      x_1=13x_3-5x_4,\\
      x_2=8x_3-3x_4,\\
      x_3\text{ is free},\\
      x_4\text{ is free}
    \end{cases}
  \]
  becomes, on setting $x_3=c_1$ and $x_4=c_2$,
  \[
    \begin{bmatrix}x_1\\x_2\\x_3\\x_4\end{bmatrix}
    =c_1\underbrace{\begin{bmatrix}13\\8\\1\\0\end{bmatrix}}_{\mathbf{u}}
     +c_2\underbrace{\begin{bmatrix}-5\\-3\\0\\1\end{bmatrix}}_{\mathbf{v}}.
  \]
\end{frame}

% ===== Slide 50 =====
\begin{frame}{Two Descriptions of a Solution Set}
  \begin{block}{Parametric description}
    A parametric description defines the coordinates of a curve, surface, or system using independent variables called parameters rather than a direct relationship between coordinates
  \end{block}
  \begin{block}{Parametric vector description}
    An explicit description of the plane as the set spanned by
    $\mathbf{u}$ and $\mathbf{v}$.
  \end{block}
\end{frame}

% ===== Slide 51 =====
\begin{frame}{Solution Sets of Homogeneous Linear Systems}
  % Chinese subtitle in source: 齐次系统
  A \alert{homogeneous system} has the form
  \[
    A\mathbf{x}=\mathbf{0},
  \]
  where $A$ is $m\times n$ and $\mathbf{0}$ is the zero vector in
  $\mathbb{R}^m$.

  \medskip
  Example:
  \[
    \begin{cases}
      x_1+10x_2=0,\\
      2x_1+20x_2=0,
    \end{cases}
    \qquad
    \begin{bmatrix}1&10\\2&20\end{bmatrix}
    \begin{bmatrix}x_1\\x_2\end{bmatrix}
    =\begin{bmatrix}0\\0\end{bmatrix}.
  \]
\end{frame}

% ===== Slide 52 =====
\begin{frame}{Trivial and Nontrivial Solutions}
  The homogeneous system $A\mathbf{x}=\mathbf{0}$ always has the
  \alert{trivial solution} (Trivial: simple, obvious, expected answer).
  \[
    \mathbf{x}=\mathbf{0}.
  \]

  Here \textbf{x} and \textbf{0} are both vectors. Basically, it is setting all vairables to 0. It always works!

  Nonzero vector solutions are called \alert{nontrivial solutions}.

  \medskip
  For the previous example,
  \[
    \left[\begin{array}{cc|c}1&10&0\\2&20&0\end{array}\right]
    \sim
    \left[\begin{array}{cc|c}1&10&0\\0&0&0\end{array}\right].
  \]
  The system is consistent and has a free variable, so it has infinitely
  many solutions, including nontrivial ones.
\end{frame}

% ===== Slide 53 =====
\begin{frame}{When Do Nontrivial Solutions Exist?}
  A homogeneous equation $A\mathbf{x}=\mathbf{0}$ has nontrivial solutions
  if and only if the system has \alert{at least one free variable}.

  \medskip
  Determine whether the following system has nontrivial solutions, and then
  describe its solution set:
  \[
    \begin{cases}
      2x_1+4x_2-6x_3=0,\\
      4x_1+8x_2-10x_3=0.
    \end{cases}
  \]
  There are three variables but at most two pivots, so at least one variable
  is free; hence nontrivial solutions exist.
\end{frame}

\begin{frame}{When Do Nontrivial Solutions Exist?}
  Why we need to have free variables for nontrivial solution to exist?
\end{frame}

% ===== Slide 54 =====
\begin{frame}{Row Reduction of the Homogeneous System}
  \[
    \left[\begin{array}{ccc|c}2&4&-6&0\\4&8&-10&0\end{array}\right]
    \sim
    \left[\begin{array}{ccc|c}1&2&-3&0\\4&8&-10&0\end{array}\right]
    \sim
    \left[\begin{array}{ccc|c}1&2&-3&0\\0&0&2&0\end{array}\right]
    \sim
    \left[\begin{array}{ccc|c}1&2&0&0\\0&0&1&0\end{array}\right].
  \]
  Therefore
  \[
    x_1=-2x_2,\qquad x_2\text{ is free},\qquad x_3=0.
  \]
\end{frame}

% ===== Slide 55 =====
\begin{frame}{Parametric Vector Description}
  Since $x_1=-2x_2$, $x_2$ is free, and $x_3=0$,
  \[
    \mathbf{x}
    =\begin{bmatrix}x_1\\x_2\\x_3\end{bmatrix}
    =\begin{bmatrix}-2x_2\\x_2\\0\end{bmatrix}
    =x_2\underbrace{\begin{bmatrix}-2\\1\\0\end{bmatrix}}_{\mathbf{v}}
    = x_2 \mathbf{v}.
  \]
  Thus the solution set is
  \[
    \operatorname{Span}\left\{\begin{bmatrix}-2\\1\\0\end{bmatrix}\right\}.
  \]
\end{frame}

% ===== Slide 56 =====
\begin{frame}{Solution Sets of Nonhomogeneous Linear Systems}
  Describe the solution set of
  \[
    \begin{cases}
      2x_1+4x_2-6x_3=0,\\
      4x_1+8x_2-10x_3=4.
    \end{cases}
  \]

  Nonhomogeneous because (0,4) is not all zero (0,0).

  Row reduction gives
  \[
    \left[\begin{array}{ccc|c}2&4&-6&0\\4&8&-10&4\end{array}\right]
    \sim
    \left[\begin{array}{ccc|c}1&2&0&6\\0&0&1&2\end{array}\right].
  \]
  Hence
  \[
    \mathbf{x}
    = \begin{bmatrix}6-2x_2\\x2\\2\end{bmatrix}
    =\underbrace{\begin{bmatrix}6\\0\\2\end{bmatrix}}_{\mathbf{p}}
     +x_2\underbrace{\begin{bmatrix}-2\\1\\0\end{bmatrix}}_{\mathbf{v}}
    = \mathbf{p}+2\mathbf{v}
  \]
\end{frame}

% ===== Slide 57 =====
\begin{frame}{Recap of the Previous Two Examples}
  The solution set of $A\mathbf{x}=\mathbf{0}$ is
  \[
    \mathbf{x}=x_2\begin{bmatrix}-2\\1\\0\end{bmatrix}=x_2\mathbf{v}.
  \]
  The solution set of $A\mathbf{x}=\mathbf{b}$ is
  \[
    \mathbf{x}
    =\begin{bmatrix}6\\0\\2\end{bmatrix}
     +x_2\begin{bmatrix}-2\\1\\0\end{bmatrix}
    =\mathbf{p}+x_2\mathbf{v}.
  \]
\end{frame}

% ===== Slide 58 =====
\begin{frame}{Theorem 6}
  Suppose $A\mathbf{x}=\mathbf{b}$ is consistent (exists solution) for some given $\mathbf{b}$ and let $\mathbf{p}$ be
  one particular solution. Then its solution set is the set of all vectors
  of the form
  \[
    \mathbf{w}=\mathbf{p}+\mathbf{v}_h,
  \]
  where $\mathbf{v}_h$ is any solution of the homogeneous equation
  $A\mathbf{x}=\mathbf{0}$.

  \medskip
  \alert{Warning:} This theorem applies only when
  $A\mathbf{x}=\mathbf{b}$ has at least one solution. If the $A\mathbf{x}=\mathbf{b}$ is 
  inconsistent, its solution set is empty.

  {\color{red}{\textbf{CHECK IF SOLUTION EXISTS FIRST!}}}

  % Corrected from the source slide: the relevant condition is that a
  % particular solution exists; it need not be described as "nonzero" in general.
\end{frame}

% ===== Slide 59 =====
\begin{frame}{Example: A Homogeneous Plane}
  Describe the solution set of $2x_1-4x_2-4x_3=0$ and compare it with
  the solution set of $2x_1-4x_2-4x_3=6$.
\pause

  Solution: for the first homogeneous equation,
  \[
    \begin{bmatrix}2&-4&-4&0\end{bmatrix}
    \sim\begin{bmatrix}1&-2&-2&0\end{bmatrix}.
  \]
  Therefore
  \[
    \mathbf{x}
    =\begin{bmatrix}2x_2+2x_3\\x_2\\x_3\end{bmatrix}
    =x_2\underbrace{\begin{bmatrix}2\\1\\0\end{bmatrix}}_{\mathbf{u}}
     +x_3\underbrace{\begin{bmatrix}2\\0\\1\end{bmatrix}}_{\mathbf{v}}.
  \]
\end{frame}

% ===== Slide 60 =====
\begin{frame}{Example: A Parallel Nonhomogeneous Plane}
  For the second nonhomogeneous equation $2x_1-4x_2-4x_3=6$,
  \[
    \begin{bmatrix}2&-4&-4&6\end{bmatrix}
    \sim\begin{bmatrix}1&-2&-2&3\end{bmatrix}.
  \]
  Thus
  \[
    \mathbf{x}
    =\begin{bmatrix}3+2x_2+2x_3\\x_2\\x_3\end{bmatrix}
    =\underbrace{\begin{bmatrix}3\\0\\0\end{bmatrix}}_{\mathbf{p}}
     +x_2\underbrace{\begin{bmatrix}2\\1\\0\end{bmatrix}}_{\mathbf{u}}
     +x_3\underbrace{\begin{bmatrix}2\\0\\1\end{bmatrix}}_{\mathbf{v}}.
  \]
\end{frame}

% ===== Slide 61 =====
\begin{frame}{Exercise: True or False?}
  \small
  \begin{enumerate}
    \item If $A\mathbf{x}=\mathbf{0}$ has only the trivial solution, then
          $A\mathbf{x}=\mathbf{b}$ has a unique solution.\pause \alert{False} \\Has 0 or 1 solution
    \item If $A\mathbf{x}=\mathbf{0}$ has a nontrivial solution, then
          $A\mathbf{x}=\mathbf{b}$ has infinitely many solutions.\pause \alert{False}\\ Has 0 or infinitely many
    \item If $A\mathbf{x}=\mathbf{b}$ has infinitely many solutions, then
          $A\mathbf{x}=\mathbf{0}$ has a nontrivial solution.\pause \alert{True} \\Only trivial 
  \end{enumerate}
\end{frame}


\begin{frame}{Exercise: True or False?}
  \small
  \begin{enumerate}[start=4]
    \item If $A\mathbf{x}=\mathbf{b}$ has a unique solution, then
          $A\mathbf{x}=\mathbf{0}$ has a nontrivial solution.\pause \alert{False} \\ Only nontrivial
    \item Any $n\times n$ system of linear equations has at most $n$
          solutions.\pause \alert{False} \\ May be 0, 1, or infinitely many
    \item The system $A\mathbf{x}=\mathbf{0}$ has the trivial solution if
          and only if there is no free variable. \pause\alert{False}\\ Completely reversed. If and only if there is at least one free variable
  \end{enumerate}
\end{frame}

    
